\documentclass[
		paper=a4,
	  twoside=false,
	 fontsize=10pt,
			  titlepage,
			  plainheadsepline,
			  plainfootsepline,
			  headsepline,
			  footsepline,
		index=totoc,	% fügt Inhaltsverzeichnis in Inhaltsangabe ein
	   listof=totoc,	% fügt das Sourcenverzeichnis in Inhaltsangabe ein
 bibliography=totoc,	% fügt Literaturverzeichnis in Inhaltsangabe ein
	 abstract=false,		%needed?
	titlepage=false
]{report}

\raggedright

%\parskip2ex plus1ex minus0.5ex
\setlength{\parindent}{0.05\textwidth}

\usepackage[
		 left=3cm,
		right=3cm,
		  top=2.5cm,
	   bottom=2.5cm,
		includeheadfoot
]{geometry}

\usepackage[
		pdftex,
		colorlinks=true,
		 linkcolor=black,
		 citecolor=black,
		 filecolor=black,
		  urlcolor=black
]{hyperref}
\usepackage[all]{hypcap}

\usepackage[english, ngerman]{babel}\selectlanguage{ngerman}
\usepackage[ngerman]{translator}
\usepackage[acronym, toc, numberedsection=false]{glossaries}

\usepackage[utf8]{inputenc} %Kodierung, statt latin1 auch utf8 möglich
%\usepackage{setspace} \onehalfspacing

\usepackage{lmodern}
%\usepackage{ae,aeguill}
\usepackage{ae}

\usepackage{amsfonts}
\usepackage{amsmath}
\usepackage[automark]{scrpage2}
\pagestyle{scrheadings}
%\setlength{\headheight}{1.5\baselineskip}
\clearscrheadfoot

\usepackage[round,numbers,square]{natbib}
\bibliographystyle{alphadin}

\usepackage{graphicx}
\usepackage{epsfig}
\setlength{\unitlength}{10mm} %sets unitlength for picture-environment

\usepackage{rotating}
\usepackage{longtable}
\usepackage{multirow}
\usepackage{afterpage}

% \usepackage{eurosym} % euro sign

\usepackage{color}
\definecolor{blue}{rgb}{0.1647,0,1}
\definecolor{darkred}{rgb}{0.498,0,0.33333}
\definecolor{green}{rgb}{0.247,0.498,0.3725}
\definecolor{lightgray}{rgb}{0.95,0.95,0.95}

\usepackage{listings}
\lstset{
		language=Java,									% the language for the code
	escapeinside={(*@}{@*)},							% the comments in code that are not printed
	  basicstyle=\scriptsize\ttfamily\color{black},		% the style of the code font
	keywordstyle=\bfseries\color{darkred},				% the style for keywords
	 stringstyle=\color{blue},							% the style for strings
	commentstyle=\color{green},							% the style for comments
 identifierstyle=\color{black},							% the style for identifiers
		 numbers=left,									% where to put the line numbers
	 numberstyle=\tiny,									% the style of line numbers
showstringspaces=false,									% underline spaces in strings
	  stepnumber=1,										% the step between two numbered lines
	   numbersep=10pt,									% space between numbers and code
	  captionpos=b,										% where the caption should be printed (bottom)
		   frame=single,  								% if a frame should be printed around the code
		 tabsize=4,										% how many spaces one tab contains
       aboveskip=1cm,									% the space above the displayed listing
 backgroundcolor=\color{lightgray},							% background color to fill the listings with
			breaklines
}
% \renewcommand{\lstlistlistingname}{Quellcodeverzeichnis}
%\renewcommand{\baselinestretch}{1.5}\normalsize

% Schusterjungen/Hurenkinder ausschließen
\clubpenalty = 10000
\widowpenalty = 10000
\displaywidowpenalty = 10000

\newcommand{\figref}[1]{Abbildung~\ref{fig:#1}}
\newcommand{\secref}[1]{Abschnitt~\ref{sec:#1}}
\newcommand{\tabref}[1]{Tabelle~\ref{tbl:#1}}
\newcommand{\lstref}[1]{Listing~\ref{lst:#1}}
\newcommand{\srcref}[1]{Zeile~\ref{src:#1}}
\newcommand{\chpref}[1]{Kapitel~\ref{chp:#1}}
\newcommand{\icon}{icon~Systemhaus~GmbH\ }

% \usepackage{ifthen}
% \newcommand{\vgl}[2][]{% #1 = optional, #2 = notwendig
%   \ifthenelse{\equal{#1}{}}{
%        % Autor hat kein oder leeres optionales Argument angegeben
% 		(Vgl. \cite{#2})
%   }{%
%        % Autor hat optionales Argument angegeben
% 		(Vgl. \cite[#1]{#2})
%   }
% }
% \newcommand{\vglausf}[2][]{% #1 = optional, #2 = notwendig
%   \ifthenelse{\equal{#1}{}}{
%        % Autor hat kein oder leeres optionales Argument angegeben
% 		(Vgl. dazu ausführlich \cite{#2})
%   }{%
%        % Autor hat optionales Argument angegeben
% 		(Vgl. dazu ausführlich \cite[#1]{#2})
%   }
% }

\newcommand{\code}[1]{\texttt{#1}}
\newcommand{\notiz}[1]{\marginpar{\tiny#1}}
\newcommand{\glentry}[3]{\newglossaryentry{#1}{name={#2},description={{#3}}}}

\usepackage[texcoord]{eso-pic}
%\newcommand{\startprintinglogo}{
%	\AddToShipoutPicture{
%		\parbox[t]{\paperwidth}{
%			\vspace{0.5cm}
%			\hspace{16.2cm}
%			\includegraphics[width=5.36cm,height=2.0cm]{includes/dhbw}
%		}
%	}
%}
